\usepackage[spanish,es-nodecimaldot]{babel}
\usepackage{iftex}
\ifPDFTeX
  \usepackage[T1]{fontenc}
  \usepackage[utf8]{inputenc}
  \usepackage{lmodern}
\else
  \usepackage{fontspec}
  \setmainfont{Latin Modern Roman}
  \setsansfont{Latin Modern Sans}
  \setmonofont{Latin Modern Mono}
\fi
\usepackage[a4paper,margin=2.4cm]{geometry}
\usepackage{amsmath,amssymb,mathtools}
\usepackage{siunitx}
\usepackage{enumitem}
\usepackage[most]{tcolorbox}
\usepackage{fancyhdr}
\usepackage{hyperref}
\usepackage{booktabs}
\usepackage{tabularx}
\usepackage{multicol}
\usepackage{xcolor}
\usepackage{tikz}
\usepackage{pgfplots}
\usetikzlibrary{babel}

\pgfplotsset{compat=1.18}
\hypersetup{
  colorlinks=true,
  linkcolor=black,
  urlcolor=black,
  pdftitle={Cuaderno de repaso de Matemáticas de 1.º de Bachillerato},
  pdfauthor={Codex}
}

\setlength{\parindent}{0pt}
\setlength{\parskip}{0.5em}
\setlength{\headheight}{14pt}
\raggedbottom
\setlist[itemize]{topsep=0.3em,itemsep=0.2em}
\setlist[enumerate]{topsep=0.3em,itemsep=0.2em}

\pagestyle{fancy}
\fancyhf{}
\fancyhead[L]{Cuaderno de repaso de Matemáticas}
\fancyhead[R]{\ifteacheredition Edición del profesorado\else Edición del alumnado\fi}
\fancyfoot[C]{\thepage}

\newcommand{\sen}{\operatorname{sen}}
\newcommand{\tg}{\operatorname{tg}}
\newcommand{\cotg}{\operatorname{cotg}}
\newcommand{\R}{\mathbb{R}}
